% Tabla: Validación del peso de muestras sintéticas
% Experimento 07a - K-Fold CV (5×3 = 15 folds)
\begin{table}[htbp]
\centering
\caption{Impacto del peso de muestras sintéticas en el entrenamiento}
\label{tab:weight_validation}
\begin{tabular}{cccccc}
\toprule
\textbf{Peso ($w$)} & \textbf{N Sintéticos} & \textbf{Calidad} & \textbf{Macro F1} & \textbf{$\Delta$ F1 (pp)} & \textbf{p-valor} \\
\midrule
0.3 & 125 & 0.231 & 0.2103 & +0.58 & $<$0.0001 \\
0.5 & 125 & 0.231 & 0.2081 & +0.36 & 0.0009 \\
0.7 & 125 & 0.231 & 0.2100 & +0.55 & $<$0.0001 \\
\rowcolor{green!20}
\textbf{1.0} & \textbf{125} & \textbf{0.231} & \textbf{0.2110} & \textbf{+0.65} & $<$\textbf{0.0001} \\
\bottomrule
\end{tabular}
\vspace{0.5em}\footnotesize

\textbullet\ Línea base F1: 0.2045. Todas las configuraciones son estadísticamente significativas.

\textbullet\ \textbf{Hallazgo clave}: peso igual ($w$=1.0) supera significativamente al peso reducido ($w$=0.5).

\textbullet\ Esto contradice la suposición inicial de que sintéticos deberían ponderarse menos que reales.

\end{table}
